\section{SF6 revision 2: fixed-gain current-state fusion}
\label{sec:sf6}
SF6's state order is
\begin{equation}
\mat x_s=[v_E,v_N,v_U,p_E,p_N,p_U]^T.
\end{equation}
It uses the admitted inertial navigation increment to predict
\begin{align}
\mat v^-_{k+1}&=\mat v^+_k+\Delta\mat v_n,\\
\mat p^-_{k+1}&=\mat p^+_k+(\mat v^+_k+\tfrac12\Delta\mat v_n)h.
\end{align}
The separate attitude propagates inertially and is not corrected by these observations. For each observed state component $j$, the update is
\begin{equation}
x_{s,j}^{+}=x_{s,j}^{-}+k_j(z_j-x_{s,j}^{-}),\qquad0\le k_j\le1.
\label{eq:sf6-gain}
\end{equation}
The six gains follow the state order above; their declared defaults are all 0.2. This is an engineering default, not an optimal gain or accuracy result. $k_j=0$ retains the prediction and $k_j=1$ selects the admitted observation. SF6 maintains no $P$, $Q$, $R$, Kalman gain or NIS. The shared output's statistical fields are unavailable (NaN), not zero uncertainty.

\begin{center}
\begin{tabular}{@{}cll@{}}
\toprule Group & Observation & Corrected components \\
\midrule
0 & GNSS position EN & $p_E,p_N$ \\
1 & GNSS position U & $p_U$ \\
2 & GNSS velocity EN & $v_E,v_N$ \\
3 & GNSS velocity U & $v_U$ \\
4 & Relative barometer altitude & $p_U$ \\
\bottomrule
\end{tabular}
\end{center}
GNSS altitude and barometer altitude share the $p_U$ gain. Position observations do not directly correct velocity. There is no magnetic measurement update in SF6.

Revision 2 applies every admitted observation to the \emph{present} predicted state. It has no history ring, delayed observation queue, rewind or replay. Measurement time serves admission/age checks. The backend uses a trusted sample timestamp, otherwise the receive timestamp (zero configured delay). A measurement must satisfy
\begin{equation}
t_0\le t_s\le t_r\le t_e,\qquad t_e-t_s\le0.5\ \mathrm{s}.
\end{equation}
For a previously seen group, both sample and receive times must advance. Unsigned 32-bit sequence difference must lie in $[1,2^{31}-1]$, which admits wraparound only within the forward half range. Prediction requires $0<h\le0.02$ s and exact continuity with the microsecond-rounded duration.

A constant-velocity delayed position observation lags current position by about $-\mat v\tau$. Equation~\eqref{eq:sf6-gain} therefore adds approximately $-k_jv_j\tau$ beyond a current observation's correction. Under constant acceleration, delayed velocity analogously contributes $-k_ja_j\tau$. These are explanatory model estimates, not measured navigation accuracy.

The backend checks source identity, basic physical GNSS/barometer validity and age. SF6 does not obtain covariance-based satellite/window/NIS downweighting simply because those policies exist for KF6/ESKF15. Its health reports admission and whether the relevant gain can make an effective update; a zero gain cannot refresh successful-fusion health. The covariance-free core commits the finite candidate state and group admission metadata together.

Historical SF6 revision 1 used fixed-lag replay. Recorded mode must retain that revision's semantics; current Offline What-if revision 2 must identify its semantic difference. Continuing inertial propagation or barometer correction during GNSS loss does not bound horizontal drift.

\source{S3}{fixed gains, state advancement, group/time admission and candidate commit}
\source{S4}{physical admission, current-state measurement timestamps and health/statistical output}
